\section{Boîtier}
\label{sec:boitier}

Le boîtier tient la carte et la batterie, il protège l'électronique et il met à disposition de l'utilisateur les commandes et les connecteurs. La conception mécanique ne fait pas partie du cadre de la formation, et le boîtier a donc été traité comme un support de l'électronique plutôt que comme un objet de conception à part entière. Les solutions retenues sont volontairement simples et rapides à mettre en œuvre, l'objectif étant de disposer sans délai d'un appareil manipulable pour la mise au point du firmware. Les plans cotés de chaque pièce figurent en annexe~\ref{ann:plans} et font foi en cas d'écart avec le texte.

\subsection{Démarche et contraintes}

Le boîtier a été dessiné après la carte, comme annoncé à la section~\ref{sec:layout}. Le circuit imprimé terminé a été exporté depuis Altium au format \gls{step}, puis importé dans Autodesk Inventor où il constitue la pièce de référence de l'assemblage. Ce modèle apporte les dimensions réelles de la carte, la hauteur de chaque composant et l'emplacement exact des connecteurs, de l'écran et des boutons, et toute la géométrie intérieure en découle. L'encombrement, la position de chaque ouverture, le logement de la batterie et le dégagement à réserver au-dessus des composants les plus hauts sont donc entièrement fixés par la carte. La seule liberté réelle porte sur la forme extérieure, reprise de la disposition d'une Game Boy~\cite{wikipedia_gameboy,copetti_gameboy_architecture}, avec deux boutons ronds et une croix directionnelle. Cette disposition a été retenue parce qu'elle est simple, immédiatement comprise par n'importe quel utilisateur sans apprentissage préalable, et directement transposable en pièces imprimables sans mécanisme particulier. Il s'agit d'un choix de confort, aucune contrainte technique ne l'impose.

\subsection{Découpage en pièces}

Le boîtier compte six pièces, représentées à la figure~\ref{boitier_eclate.png}. Une coque avant et une coque arrière enferment la carte, un couvercle ferme le logement de la batterie sur la face arrière, et trois boutons prennent place dans les ouvertures de la face avant.

\fig[H, width=0.75\linewidth]{Vue éclatée du boîtier, coque avant, boutons, circuit, coque arrière et couvercle de batterie}{boitier_eclate.png}

Ce découpage répond à trois besoins. Le plan de séparation entre les deux coques est parallèle à la carte, ce qui permet d'insérer celle-ci par le dessus et donne à chaque coque une grande face plane à poser sur le plateau d'impression. Le couvercle de batterie est une pièce distincte afin que la cellule puisse être remplacée sans ouvrir l'appareil. Les trois boutons sont séparés parce qu'ils doivent se déplacer, leur rôle étant de transmettre l'appui du doigt aux boutons soudés sur la face avant de la carte.

\subsection{Ouvertures, fixation et fabrication}

La fenêtre d'écran, les trous des boutons, la prise casque, le connecteur USB-C, le lecteur de carte microSD et le potentiomètre de volume sont définis en fonction des composants du modèle importé, dont la position sur la carte a été fixée lors du layout.

L'assemblage repose sur de la visserie M3. Un filetage imprimé directement dans le plastique tient mal et se dégrade au démontage, la solution retenue consiste donc à noyer quatre écrous dans la coque avant. Un logement à leurs dimensions est prévu dans la pièce, l'impression est mise en pause lorsque ce logement arrive à hauteur, les écrous y sont déposés, puis l'impression reprend et referme la matière au-dessus d'eux. Les écrous ne peuvent alors ni tourner ni ressortir, et l'assemblage se fait ensuite par de simples vis traversantes, sans insert à sertir à chaud ni outillage particulier.

La carte n'est vissée sur aucune des deux coques. Elle repose sur des colonnettes plates de la coque avant, qui fixent sa hauteur à l'intérieur du boîtier, et la coque arrière vient la plaquer contre ces appuis. Quatre vis introduites par la coque arrière se serrent dans les écrous noyés, de sorte qu'un même serrage assure la fermeture du boîtier et le maintien du circuit imprimé. Le couvercle de batterie est retenu à part par une vis unique, ce qui permet de le retirer sans desserrer les quatre autres.

Les boutons sont laissés libres dans leurs logements. Chacun porte sur sa face intérieure un épaulement plus large que l'ouverture, qui l'empêche de sortir du boîtier une fois la coque avant en place. Aucun guidage ni ressort de rappel n'a été ajouté, la course étant limitée par cet épaulement d'un côté et par le bouton de la carte de l'autre, ce dernier assurant seul le retour en position haute. Les deux boutons ronds restent libres en rotation, sans effet sur leur fonction, tandis que la croix directionnelle est retenue par la forme de son ouverture, qui reprend son contour.

Les six pièces sont imprimées en \gls{fdm}. Le découpage et l'orientation d'impression ont été choisis pour que chaque pièce repose sur une grande face plane et ne demande pas de support à l'intérieur des logements.

\subsection{Validation du boîtier}

La première vérification est virtuelle. Le modèle \gls{step} de la carte étant présent dans l'assemblage, l'analyse d'interférences de l'outil de \gls{cao} a été menée sur l'ensemble des pièces afin de contrôler qu'aucun composant ne touche une paroi et que le dégagement au-dessus des plus hauts d'entre eux est suffisant. Le même modèle a servi à contrôler ouverture par ouverture que chaque percement se trouve bien en face du connecteur qu'il dessert.

La seconde vérification est physique, le boîtier ayant été imprimé et monté. Le montage confirme que les six pièces s'assemblent, que les écrous noyés reprennent le serrage des quatre vis, que la carte est correctement plaquée sur ses colonnettes et que les connecteurs restent accessibles depuis l'extérieur. Les trois boutons ont été contrôlés un à un jusqu'à l'actionnement effectif du contact correspondant sur la carte, et le retrait du couvercle donne accès à la batterie sans démontage du reste de l'appareil.